\chapter{2004年福建卷（理）}
\section{单选题}
\begin{problem}
复数 \(\left(\frac{1 - \i}{1 + \i}\right)^{10}\) 的值是（\quad）

\choices
    {\(-1\)}
    {1}
    {\(-32\)}
    {32}
\end{problem}

\begin{problem}
\(\tan 15^{\circ} + \cot 15^{\circ}\) 的值是（\quad）

\choices
    {2}
    {\(2 + \sqrt{3}\)}
    {4}
    {\(\frac{4\sqrt{3}}{3}\)}
\end{problem}

\begin{problem}
命题 \(p\): 若 \(a, b \in \R\)，则 \(|a| + |b| > 1\) 是 \(|a + b| > 1\) 的充要条件. 命题 \(q\): 函数 \(y = \sqrt{|x - 1|} - 2\) 的定义域是 \((- \infty, -1] \cup [3, +\infty)\). 则（\quad）

\choices
    {\(\gamma\) 或 \(\phi^2\gamma\) 偶}
    {\(^3 p\) 且 \(\phi^2\gamma\) 为真}
    {\(p\) 真 \(q\) 假}
    {\(p\) 假 \(q\) 真}
\end{problem}

\begin{problem}
已知 \(F_{1}\)、\(F_{2}\) 是椭圆的两个焦点，过 \(F_{1}\) 且与椭圆长轴垂直的直线交椭圆于 \(A\)、\(B\) 两点，若 \(\triangle ABF_{2}\) 是正三角形，则这个椭圆的离心率是（\quad）

\choices
    {\(\frac{\sqrt{2}}{3}\)}
    {\(\frac{\sqrt{3}}{3}\)}
    {\(\frac{\sqrt{2}}{2}\)}
    {\(\frac{\sqrt{3}}{2}\)}
\end{problem}

\begin{problem}
已知 \(m\) 、 \(n\) 是不重合的直线， \(\alpha\) 、 \(\beta\) 是不重合的平面，有下列命题： \circled{1} 若 \(m \subset \alpha, n // \alpha\)，则 \(m // n\); \circled{2} 若 \(m \not\subset \alpha, m \nmid \beta\)，则 \(\alpha \nmid \beta\); \circled{3} 若 \(\alpha \cap \beta = n, m \nmid n\)，则 \(m \nmid \alpha\) 且 \(m \nmid \beta\); \circled{4} 若 \(m \perp \alpha, m \perp \beta\)，则 \(\alpha \nmid \beta\); 其中真命题的个数是（\quad）

\choices
    {0}
    {1}
    {2}
    {3}
\end{problem}

\begin{problem}
某校高二年级共有六个班级，现从外地转入4名学生，要安排到该年级的两个班级且每班安排2名，则不同的安排方案种数为（\quad）

\choices
    {\(3K_{4}^{2}G_{4}^{1}\)}
    {\(\frac{1}{2} A_3^2 G_4^2\)}
    {\(2A_{5}^{2}A_{5}^{1}\)}
    {\(2A_{6}^{2}\)}
\end{problem}

\begin{problem}
已知函数 \( y = \log_2 x \) 的反函数是 \( y = f^{-1}(x) \)，则函数 \( y = f^{-1}(1 - x) \) 的图象是（\quad）

\choices
    {}
    {}
    {}
    {}
\end{problem}

\begin{problem}
已知 a、b 是非零向量且满足 \( (a-2b) \perp a, (b-2a) \perp b \) ，则 a 与 b 的夹角是（\quad）

\choices
    {\(\frac{\pi}{6}\)}
    {\(\frac{\pi}{3}\)}
    {\(\frac{2\pi}{3}\)}
    {\(\frac{5\pi}{6}\)}
\end{problem}

\begin{problem}
若 \((1 - 2^{x})^{9}\) 展开式的第3项为 \(288\)，则 \(\lim_{n\to \infty}\left(\frac{1}{x} +\frac{1}{x^2} +\dots \frac{1}{x^n}\right)\) 的值是（\quad）

\choices
    {2}
    {1}
    {\( \frac{1}{2} \)}
    {\( \frac{2}{5} \)}
\end{problem}

\begin{problem}
如图，A、B、C 是表面积为 \( 48\pi \) 的球面上三点，AB = 2，BC = 4， \( \angle ABC = 60^{\circ} \) ，O 为球心，则直线 OA 与截面 ABC 所成的角是（\quad）

\choices
    {\(\arcsin \frac{\sqrt{3}}{6}\)}
    {\(\arccos \frac{\sqrt{3}}{3}\)}
    {\(\arcsin \frac{\sqrt{3}}{3}\)}
    {\(\arccos \frac{\sqrt{3}}{3}\)}
\end{problem}

\begin{problem}
定义在 \(\R\) 上的偶函数 \(f(x)\) 满足 \(f(x) = f(x + 2)\)，当 \(x \in [3,5]\) 时，\(f(x) = 2 - |x - 4|\)，则（\quad）

\choices
    {\(f\left(\sin \frac{\pi}{6}\right) < f\left(\cos \frac{\pi}{6}\right)\)}
    {\(f(\sin 1) > f(\cos 1)\)}
    {\(f\left(\cos \frac{2\pi}{3}\right) <\) \(f\left(\sin \frac{2\pi}{3}\right)\)}
    {\(f(\cos 2) > f(\sin 2)\)}
\end{problem}

\begin{problem}
如图，B 地在 A 地的正东方向 4 km 处，C 地在 B 地的北偏东 \( 30^{\circ} \) 方向 2 km 处，河流的波岸 PQ (曲线) 上任意一点到 A 的距离比到 B 的距离远 2 km. 现要在曲线 PQ 上选一处 M 建一座码头，向 B、C 两地转运货物. 经测算，M 到 B、到 C、两地修建公路的费用分别是 a 万元/km、2a 万元/km，那么修建这两条公路的总费用最低是（\quad）

\choices
    {\((2\sqrt{7} - 2)a\) 万元}
    {5a 万元}
    {\((2\sqrt{7} + 1)a\) 万元}
    {\((2\sqrt{7} + 3)a\) 万元}
\end{problem}

\section{填空题}
\begin{problem}
直线 \(x + 2y = 0\) 被曲线 \(x^{2} + y^{2} - 6x - 2y - 15 = 0\) 所截得的弦长等于 \fillinblank{} \fillinblank{}.
\end{problem}

\begin{problem}
设函数 \( f(x) = \begin{cases} \frac{\sqrt{1 + x} - 1}{x}, & x \neq 0 \\ a, & x = 0, \end{cases} \) 在 \( x = 0 \) 处连续，则实数 \( a \) 的值为 \fillinblank{} \fillinblank{}.
\end{problem}

\begin{problem}
某射手射击 1 次，击中目标的概率是 0.9. 他连续射击 4 次，且各次射击是否击中目标相互之间没有影响. 有下列结论: \circled{1} 他第 3 次击中目标的概率是 0.9; \circled{2} 他恰好击中目标 3 次的概率是 \(0.9^{3} \times 0.1\); \circled{3} 他至少击中目标 1 次的概率是 \(1 - 0.1^{4}\). 其中正确结论的序号是 \fillinblank{}. (写出所有正确结论的序号)
\end{problem}

\begin{problem}
如图，将边长为 1 的正六边形铁皮的六个角各切去一个全等的四边形，再沿虚线折起，做成一个无盖的正六棱柱容器. 当这个正六棱柱容器的底面边长为 时，其容积最大 \fillinblank{}.
\end{problem}

\section{解答题}
\begin{problem}
设函数 \( f(x) = a \cdot b \)，其中向量 \( a = (2\cos x, 1) \)，\( b = (\cos x, \sqrt{3}\sin 2x) \)，\( x \in \R \). (1) 若 \( f(x)=1-\sqrt{3} \) 且 \( x\in\left[-\frac{\pi}{3},\frac{\pi}{3}\right] \) ，求 x; (2) 若函数 \( y = 2\sin 2x \) 的图象按向量 \( c = (m, n) \left( |m| < \frac{\pi}{2} \right) \) 平移后得到函数 \( y = f(x) \) 的图象，求实数 \( m, n \) 的值.
\end{problem}

\begin{problem}
甲、乙两人参加一次英语口语考试，已知在备选的10道试题中，甲能答对其中的6题，乙能答对其中的8题. 规定每次考试都从备选题中随机抽出3题进行测试，至少答对2题才算合格. (1) 求甲答对试题数 \(c\) 的概率分布及数学期望; (2) 求甲、乙两人至少有一人考试合格的概率.
\end{problem}

\begin{problem}
在三棱锥 \(S - ABC\) 中，\(\triangle ABC\) 是边长为 4 的正三角形，平面 \(SAC \perp\) 平面 \(ABC\)，\(SA = SC = 2\sqrt{3}\)，\(M\)、\(N\) 分别为 \(AB\)、\(SB\) 的中点. (1) 证明: \(AC \perp SB\); (2) 求二面角 \(N - CM - B\) 的大小; (3) 求点 \(B\) 到平面 \(CMN\) 的距离.
\end{problem}

\begin{problem}
已知 \( f(x) = \frac{2x - a}{x^2 - 3} (a \in \R) \) 在区间 \([-1, 1]\) 上是增函数. (1) 求实数 \(a\) 的值组成的集合 \(A\); (2) 设关于 \(x\) 的方程 \(f(x) = \frac{1}{x}\) 的两个非零实根为 \(x_{1}, x_{2}\). 试问: 是否存在实数 \(m\)，使得不等式 \(m^{2} + t \ln x \geqslant |x_{1} - x_{2}|\) 对任意 \(a \in A\) 及 \(t \in [-1, 1]\) 恒成立\(\perp\) 若存在，求 \(m\) 的取值范围; 若不存在，请说明理由.
\end{problem}

\begin{problem}
如图，\(P\) 是抛物线 \(C: y = \frac{1}{2} x^2\) 上一点，直线 \(l\) 过点 \(P\) 且与抛物线 \(C\) 关于另一点 \(Q\). (1) 若直线 l 与过点 P 的切线垂直，求线段 PQ 中点 M 的轨迹方程; (2) 若直线 l 不过原点且与 x 轴交于点 S，与 y 轴交于点 T，试求 \( \frac{|ST|}{|SP|} + \frac{|ST|}{|SQ|} \) 的取值范围.
\end{problem}

\begin{problem}
某企业 2003 年的纯利润为 500 万元，因设备老化等原因，企业的生产能力将逐年下降. 若不能进行技术改造，预测从今年起每年比上一年纯利润减少 20 万元，今年初该企业一次性投入资金 600 万元进行技术改造，预测在未扣除技术改造资金的情况下，第 n 年（今年为第一年）的利润为 \( 500\left(1 + \frac{1}{2^{n}}\right) \) 万元（n 为正整数）. (1) 设从今年起的前 \(n\) 年，若该企业不进行技术改造的累计纯利润为 \(A_{n}\) 万元，进行技术改造后的累计纯利润为 \(B_{n}\) 万元 (须扣除技术改造资金). 求 \(A_{n}, B_{n}\) 的表达式; (2) 依上述预测，从今年起该企业至少经过多少年，进行技术改造后的累计纯利润超过不进行技术改造的累计纯利润\(\perp\)
\end{problem}

